% !TEX root = ../main.tex
% ---------------------------------------------------------------------
%  B6 -- Tasks, Echtzeit, EtherCAT (Skript Kap. 7; AS 6)
% ---------------------------------------------------------------------
\vlsec{cB6}{B6}{Tasks, Echtzeit, EtherCAT}
\label{sec:tasks}

\SW{Task} Ausführungseinheit: ruft die zugeordneten Programme mit fester
\T{Zykluszeit} und \T{Priorität} auf. Mehrere Tasks \ra verschiedene
Programme mit verschiedenen Zykluszeiten im selben Projekt.
\SW{Determinismus} Programm läuft exakt alle $T$. Schwankung $=$
\T{Jitter} \ra Regler stimmt nicht mehr (Entwurf setzt konstantes $T$
voraus), Roboter positioniert falsch, Druck an falscher Stelle.

\begin{formel}[Prioritäten in TwinCAT]
\T{Präemptiv}: Wird ein höher priorer Task fällig, wird der laufende
niedriger priore unterbrochen und später fortgesetzt.\\
\T{1 $=$ höchste} Priorität. \T{Jede Prioritätsstufe nur für genau einen
Task} -- zwei Tasks mit gleicher Priorität geht nicht.\\
Zykluszeit $=$ ganzzahliges Vielfaches der \T{Basiszeit} (SYSTEM \ra
Echtzeit), eingestellt in Zyklusticks.
\end{formel}
\begin{falle}[Zykluszeit und CPU-Last]
Kürzere Zykluszeit \ra Programm läuft öfter \ra \emph{höhere}
CPU-Auslastung. Last je Task $=\frac{\text{Laufzeit}}{\text{Zykluszeit}}$.
Summe $>100\,\%$ \ra Überlast, niedrig priore Tasks verpassen ihren Zyklus.
\end{falle}

\begin{bsp}[AS 6: drei Tasks, je 30\,$\mu$s Rechenzeit]
\label{bsp:tasks}
FastTask 300\,$\mu$s Prio 2, Standard 1\,ms Prio 1, SlowTask 10\,ms Prio 3.\\
\T{Basiszeit}: muss 300\,$\mu$s und 1\,ms teilen \ra $\text{ggT}=100\,\mu$s.\\
\T{Last}: $\frac{30}{300}+\frac{30}{1000}+\frac{30}{10000}=10+3+0.3=13.3\,\%$.\\
\T{$t=0$} (alle fällig): Standard 0--30, Fast 30--60, Slow 60--90\,$\mu$s.\\
\T{$t=3$\,ms} (Fast + Standard): Standard zuerst, Fast 30\,$\mu$s zu spät
\ra Jitter für Fast. Sonst laufen die Tasks allein.\\
\T{Lehre:} Der schnellste Task sollte die höchste Priorität haben, sonst
wird gerade er verzögert.
\end{bsp}

\subsection{EtherCAT}
\label{sec:ethercat}
\begin{formel}[Funktionsprinzip]
Normales Ethernet-Frame, \T{Ethertype \R{0x88A4}}. Der Master (SPS)
sendet ein Frame mit allen Ausgangsdaten; es läuft \emph{seriell} durch alle
Klemmen, die letzte schickt es zurück (Ring). Jede Klemme entnimmt ihre
Ausgänge und setzt ihre Eingänge ein.\\
\T{On-the-fly}: Verarbeitung \emph{während} das Frame durchläuft
(ns Verzögerung) statt empfangen-verarbeiten-senden.\\
\T{In-place}: Eingangsdaten überschreiben bereits gelesene Ausgangsdaten
\ra Frame bleibt kurz.\\
Master-Netzwerkkarte schreibt per \T{DMA} direkt ins RAM \ra
Protokollbearbeitung in Hardware, unabhängig von CPU-Last.
\end{formel}
\SW{Zahlen} Zyklen $\le100\,\mu$s; 256 verteilte DIO in 10\,$\mu$s,
1000 DIO in 30\,$\mu$s; Synchronisation (Distributed Clocks) Jitter
$\le1\,\mu$s; offener Standard (ETG), günstige Standard-Hardware.
\SW{Frames} Pro Task mit E/A ein eigenes Frame/Zyklus (Reiter EtherCAT
am Gerät, Spalte Cycle).
\SW{EL6601} Standard-Ethernet in den Lücken zwischen EtherCAT-Frames.
\SW{XFC / Oversampling} z.\,B.\ EL1262: $n$ Abtastungen pro Buszyklus
($1\,\text{ms}/n=10$ \ra 100\,$\mu$s), gesammelt übertragen.

\begin{ztab}{@{}lX@{}}
Begriff & Bedeutung \\ \midrule
Working Counter (WC) & Jede Klemme, die ihr Datagramm bearbeitet,
  erhöht den WC. Master vergleicht mit Sollwert: Abweichung \ra Daten
  dieser Gruppe ungültig, Diagnose. \\
Sync Unit & Gruppe von Klemmen mit eigenem WC. Fällt eine Gruppe aus,
  bleiben die Daten der anderen gültig. Konfiguration am EtherCAT-Gerät
  (Sync-Unit-Zuordnung). \\
Hot Connect & Gerätegruppen, die im Betrieb an- und abgesteckt werden
  dürfen (Werkzeugwechsler, modulare Maschinenteile) -- typischer Fall,
  wo Sync Units unabdingbar sind. \\
\end{ztab}
